\documentclass[12pt]{amsart}
% Package Header

\usepackage[english]{babel}
\usepackage{amsfonts}
\usepackage{amsmath}
\usepackage{amssymb}
\usepackage{animate}
\usepackage{array}
\usepackage{arrayjobx}
\usepackage{arydshln}
\usepackage{calculator}
\usepackage{cancel}
\usepackage{enumerate}
\usepackage{mathabx}
\usepackage{mathtools}
\usepackage{multicol}
\usepackage{multimedia}
\usepackage{multirow}
\usepackage{palatino}
\usepackage{polynom}
\usepackage{soul}
\usepackage{tcolorbox}
\usepackage{tikz}
\usepackage{wrapfig}
\usepackage{xcolor}
\usetikzlibrary{arrows,shapes,trees,shapes.geometric,calc}
\usepackage{array}
\usepackage{geometry}
\usepackage{enumitem}
\usepackage{parskip}
\usepackage{booktabs}
\usepackage{graphicx}
\usepackage{mathrsfs}

% \geometry{margin=1in, top=0.75in}
% \topmargin = -1.0in
% \textwidth = 6.6in
% \oddsidemargin = -0.9in
% \evensidemargin = -0.9in

\renewcommand{\headrulewidth}{0pt} % remove lines as well

\numberwithin{equation}{section}

\DeclareGraphicsRule{.tif}{png}{.png}{`convert #1 `dirname #1`/`basename #1 .tif`.png}

\textheight = 10.8 in
\topmargin = -1.0in
\textwidth = 7.9in
\oddsidemargin = -0.8in
\evensidemargin = -0.8in

\begin{document}

\newcommand{\NP}{\newpage\thispagestyle{empty}}
\newcommand{\B}{\textcolor{blue}}
\newcommand{\R}{\textcolor{red}}
\newcommand{\G}{\textcolor{olive}}
\newcommand{\Or}{\textcolor{YellowOrange}}
\newcommand{\DEF}{\underline{Defn}:\ \ \ }
\newcommand{\THM}{\underline{Thm}:\ \ \ }
\newcommand{\PROP}{\underline{Properties}:\ \ \ }
\newcommand{\NOTE}{\underline{Notes}:\ \ \ }
\newcommand{\EX}{\underline{Ex}:\ \ \ }
\newcommand{\SOL}{\underline{Soln}:\ \ \ }
\newcommand{\UL}{\underline}
\newcommand{\RE}{\mathbf{R}}
\newcommand{\V}{\textbf{ v }}
\newcommand{\VV}{[v_1,v_2]}
\newcommand{\VVV}{[v_1,v_2,v_3]}
\newcommand{\W}{\textbf{ w }}
\newcommand{\WW}{[w_1,w_2]}
\newcommand{\WWW}{[w_1,w_2,w_3]}
\newcommand{\VA}{\textbf{ a }}
%\newcommand{\AA}{[a_1,a_2]}
\newcommand{\AAA}{[a_1,a_2,a_3]}
\newcommand{\VB}{\textbf{ b }}
\newcommand{\BB}{[b_1,b_2]}
\newcommand{\BBB}{[b_1,b_2,b_3]}
\newcommand{\VC}{\textbf{ c }}
\newcommand{\CC}{[c_1,c_2]}
\newcommand{\CCC}{[c_1,c_2,c_3]}
\newcommand{\VI}{\textbf{ i }}
\newcommand{\VJ}{\textbf{ j }}
\newcommand{\VK}{\textbf{ k }}
\newcommand{\VF}{\textbf{ F }}
\newcommand{\VG}{\textbf{ G }}
\newcommand{\VO}{{\vec{0}}}
\newcommand{\OBUL}{\quad $\circ$ \ \ \ }
\newcommand{\NEWLINE}{\\ \hspace{0.63cm}}
\newcommand{\CAN}{\cancel}
\newcommand{\GF}{{\nabla f}}
\newcommand{\PFPX}{{\dfrac{\partial f}{\partial x}}}
\newcommand{\PFPY}{{\dfrac{\partial f}{\partial y}}}
\newcommand{\PFPZ}{{\dfrac{\partial f}{\partial z}}}
\newcommand{\DFDX}{{\dfrac{df}{dx}}}
\newcommand{\GG}{{\nabla g}}
\newcommand{\PGPX}{{\dfrac{\partial g}{\partial x}}}
\newcommand{\PGPY}{{\dfrac{\partial g}{\partial y}}}
\newcommand{\PGPZ}{{\dfrac{\partial g}{\partial z}}}
\newcommand{\DGDX}{{\dfrac{dg}{dx}}}
\newcommand{\PPX}{{\dfrac{\partial}{\partial x}}}
\newcommand{\PPY}{{\dfrac{\partial}{\partial y}}}
\newcommand{\PPZ}{{\dfrac{\partial}{\partial z}}}
\newcommand{\PP}{\partial}
\newcommand{\DELF}{\nabla \cdot \textbf{ F }}
\newcommand{\CURLF}{\nabla \times \textbf{ F }}
\newcommand{\DEL}{\nabla \cdot }
\newcommand{\CURL}{\nabla \times }
\newcommand{\RT}{\textbf{ r }(t)}
\newcommand{\RPT}{\textbf{ r }'(t)}
\newcommand{\GINT}{\iint \limits_{R} (\text{curl} \ \textbf{ F }) \cdot \textbf{ k } \ dA}
\newcommand{\ON}{\onslide}
\newcommand{\OB}{\overbrace}
\newcommand{\OBT}{\overbracket}
\newcommand{\OS}{\overset}
\newcommand{\UB}{\underbrace}
\newcommand{\UBT}{\underbracket}
\newcommand{\US}{\underset}
\newcommand{\D}{\displaystyle}
\newcommand{\MCol}{\multicolumn}
\newcommand{\RRE}{\mathbf{R}^2}
\newcommand{\RRRE}{\mathbf{R}^3}
\newcommand{\DXDT}{{\dfrac{dx}{dt}}}
\newcommand{\DYDT}{{\dfrac{dy}{dt}}}
\newcommand{\DZDT}{{\dfrac{dz}{dt}}}
\newcommand{\FRT}{\textbf{ F }(\textbf{ r }(t))}
\newcommand{\WORK}{ {\int \limits_{C} \textbf{ F } \cdot d\textbf{ r }} }
\newcommand{\WORKAB}{ {\int \limits_{A}^{B}\textbf{ F } \cdot d\textbf{ r }} }
\newcommand{\WORKT}{ {\int \limits_{a}^{b}\textbf{ F }(\textbf{ r }(t)) \cdot \textbf{ r }'(t) \ dt} }
\newcommand{\FLUX}{ {\iint \limits_{S}\textbf{ F } \cdot d\textbf{ S }} }
\newcommand{\FLUXX}{ {\iint \limits_{S}\textbf{ F } \cdot \textbf{ n } \ dS} }
\newcommand{\FLUXXX}{ {\iint \limits_{R}\textbf{ F }(\textbf{ r }(u,v)) \cdot \textbf{ N } \ du\ dv} }
\newcommand{\OWORK}{ {\oint \limits_{C}\textbf{ F } \cdot d\textbf{ r }} }
\newcommand{\GREENT}{ {\oint \limits_{C}\textbf{ F }\cdot \textbf{ T } \ ds} }
\newcommand{\GREENN}{ {\oint \limits_{C}\textbf{ F }\cdot \textbf{ N } \ ds} }
\newcommand{\DIVTHM}{\iiint \limits_{T} \text{div} \ \textbf{ F }  \ dV}
\newcommand{\GREENCURL}{\iint \limits_{R} (\text{curl} \ \textbf{ F }) \cdot \textbf{ k } \ dA}
\newcommand{\STOKECURL}{\iint \limits_{S} (\text{curl} \ \textbf{ F }) \cdot \textbf{ n } \ dS}
\newcommand{\GREENDIV}{\iint \limits_{R} \text{div} \ \textbf{ F } \ dA}
\newcommand{\GREENSA}{\iint \limits_{R} \left(\dfrac{\partial N}{\partial x} - \dfrac{\partial M}{\partial y}\right) \ dA}
\newcommand{\GREENSL}{\oint \limits_{C} (M\ dx + N\ dy )}
\newcommand{\FDR}{ {\textbf{ F } \cdot d\textbf{ r }} }
\newcommand{\FDRT}{ {\textbf{ F }(\textbf{ r }(t)) \cdot \textbf{ r }'(t)} }
\newcommand{\BCOLS}{\begin{columns}}
\newcommand{\BCOL}{\begin{column}}
\newcommand{\ECOLS}{\end{columns}}
\newcommand{\ECOL}{\end{column}}
\newcommand{\BLINE}{\tikz[baseline]}
\newcommand{\NODERRB}{\node[fill=red!10,draw,rectangle,anchor=base]}
\newcommand{\NODEGRB}{\node[fill=green!10,draw,rectangle,anchor=base]}
\newcommand{\NODEBRB}{\node[fill=blue!10,draw,rectangle,anchor=base]}
\newcommand{\NODECRB}{\node[fill=red!0,draw,rectangle,anchor=base]}
\newcommand{\NODERR}{\node[fill=red!10,rectangle,anchor=base]}
\newcommand{\NODEGR}{\node[fill=green!10,rectangle,anchor=base]}
\newcommand{\NODEBR}{\node[fill=blue!10,rectangle,anchor=base]}
\newcommand{\NODECR}{\node[fill=red!0,rectangle,anchor=base]}
\newcommand{\NODERE}{\node[fill=red!10,draw,ellipse,anchor=base]}
\newcommand{\NODEGE}{\node[fill=green!10,draw,ellipse,anchor=base]}
\newcommand{\NODEBE}{\node[fill=blue!10,draw,ellipse,anchor=base]}
\newcommand{\NODECE}{\node[fill=red!0,draw,ellipse,anchor=base]}
\newcommand{\NODERC}{\node[fill=red!10,draw,circle,anchor=base]}
\newcommand{\NODEGC}{\node[fill=green!10,draw,circle,anchor=base]}
\newcommand{\NODEBC}{\node[fill=blue!10,draw,circle,anchor=base]}
\newcommand{\NODECC}{\node[fill=blue!0,draw,circle,anchor=base]}
\newcommand{\CNODE}{\tikz[na]\node [coordinate] }
\newcommand{\BENDL}{edge [bend left] }
\newcommand{\BENDR}{edge [bend right] }
\newcommand{\ARROWR}{\path [red,->]}
\newcommand{\ARROWG}{\path [green,->]}
\newcommand{\ARROWB}{\path [blue,->]}
\newcommand{\ARROWSR}{\path [red,<->]}
\newcommand{\ARROWSG}{\path [green,<->]}
\newcommand{\ARROWSB}{\path [blue,<->]}
\newcommand{\LINER}{\path [red,-]}
\newcommand{\LINEG}{\path [green,--]}
\newcommand{\LINEB}{\path [blue,--]}
\newcommand{\DE}{east}
\newcommand{\DW}{west}
\newcommand{\DN}{north}
\newcommand{\DS}{south}

\newcommand{\TT}[1]{\texttt{#1}}

\tikzstyle{every picture}+=[remember picture]
\tikzstyle{na} = [baseline=-1.5ex]

%%-----------------------------------------------------------
%% Bliss Definitions
\newcommand{\To}{\Rightarrow}
\newcommand{\del}{\nabla}
\newcommand{\C}{\mathcal{C}}
\newcommand{\N}{\mathcal{N}}
\newcommand{\A}{\mathcal{A}}
\renewcommand{\O}{\mathcal{O}}
% \renewcommand{\N}{\mathbb{N}}
\newcommand{\ben}{\begin{enumerate}}
\newcommand{\een}{\end{enumerate}}
\newcommand{\eps}{\varepsilon}
\renewcommand{\epsilon}{\varepsilon}
\newcommand{\p}{\partial}
\newcommand{\<}{\langle}
\renewcommand{\>}{\rangle}
\renewcommand{\subset}{\subseteq}
\newcommand{\cont}{\Rightarrow \Leftarrow}
\newcommand{\back}{\backslash}
\newcommand{\abs}[1]{|{#1}|}
\newcommand{\ip}[1]{\langle{#1}\rangle}
\newcommand{\m}[1]{{\mathbf{#1}}}
\newcommand{\sts}{\m{S}^\textrm{T}\m{S}}
\newcommand{\rank}{\textrm{rank}}
\newcommand{\X}{\mathcal{X}}
\newcommand{\Y}{\mathcal{Y}}
\def\ds{\displaystyle}
\def\d{\partial}
\renewcommand{\arraystretch}{1.5}
%%-----------------------------------------------------------


\newcommand\CEN[1]{\multicolumn{1}{c}{ #1 }}

\newcommand\MATL[3]
{{%
\arraycolsep=#1\arraycolsep\ensuremath{{\setlength{\extrarowheight}{#2cm}\begin{bmatrix*}[l]#3\end{bmatrix*}}}
}}
\newcommand\MATC[3]
{{%
\arraycolsep=#1\arraycolsep\ensuremath{{\setlength{\extrarowheight}{#2cm}\begin{bmatrix*}[c]#3\end{bmatrix*}}}
}}
\newcommand\MATR[3]
{{%
\arraycolsep=#1\arraycolsep\ensuremath{{\setlength{\extrarowheight}{#2cm}\begin{bmatrix*}[r]#3\end{bmatrix*}}}
}}

\newcommand\MATDL[3]
{{%
\arraycolsep=#1\arraycolsep\ensuremath{{\setlength{\extrarowheight}{#2cm}\begin{vmatrix*}[l]#3\end{vmatrix*}}}
}}
\newcommand\MATDC[3]
{{%
\arraycolsep=#1\arraycolsep\ensuremath{{\setlength{\extrarowheight}{#2cm}\begin{vmatrix*}[c]#3\end{vmatrix*}}}
}}
\newcommand\MATDR[3]
{{%
\arraycolsep=#1\arraycolsep\ensuremath{{\setlength{\extrarowheight}{#2cm}\begin{vmatrix*}[r]#3\end{vmatrix*}}}
}}

\newcommand\TABwB[3]
{{
\setlength{\extrarowheight}{#1 cm}
\begin{tabular}{#2}
	\hline
	#3
	\hline
\end{tabular}
}}

\newcommand\TABwoB[3]
{{
\setlength{\extrarowheight}{#1cm}
\begin{tabular}{#2}
	#3
\end{tabular}
}}

\newcommand\includegraphicsDD[2]
{%
\includegraphics[scale=#1]{#2}
}

\newcommand\includegraphicsDDD[2]
{%
\includemedia[scale=#1,3Dmenu,activate=onclick,deactivate=onclick,
3Droll=0.,
3Dortho=0.0046491301618516445,
3Dc2c=0.8617081642150879 -0.29301872849464417 0.4142450988292694,
3Dcoo=-13.165555000305176 3.489872455596924 52.97535705566406,
3Droo=149.99999574342704,
3Dlights=Headlamp,add3Djscript=asylabels.js]{\includegraphics{#2}}
{#2.prc}
}

\newcommand\MATS[2]
{{%
  \arraycolsep=#1\arraycolsep\ensuremath{\begin{bmatrix*}[r]#2\end{bmatrix*}}
}}

\newcommand{\OKFACE}{

\begin{tikzpicture}
\draw[thick] (0cm,0cm) circle(.2cm);
\draw[thick] (-0.1,-0.09) -- (0.1,-0.09);
\draw [thick, fill=black] (0,0) circle (0.01);
\draw [rotate=90, fill=white] (0.06,0.06) circle (0.03);
\draw [rotate=90, fill=white] (0.06,-0.06) circle (0.03);
\end{tikzpicture}
}

\newcommand{\HAPPYFACE}{

\begin{tikzpicture}
\draw[thick] (0cm,0cm) circle(.2cm);
\draw[thick] plot [smooth,tension=1.5] coordinates{(-0.1,-0.08) (0,-0.13) (0.1,-0.08)};
\draw [thick, fill=black] (0,0) circle (0.01);
\draw [rotate=90, fill=white] (0.06,0.06) circle (0.04);
\draw [rotate=90, fill=white] (0.06,-0.06) circle (0.04);
\end{tikzpicture}
}

\newcommand{\SADFACE}{

\begin{tikzpicture}
\draw[thick] (0cm,0cm) circle(.2cm);
\draw[thick] plot [smooth,tension=1.5] coordinates{(-0.1,-0.12) (0,-0.07) (0.1,-0.12)};
\draw [thick, fill=black] (0,0) circle (0.01);
\draw [rotate=90, fill=white] (0.06,0.06) circle (0.02);
\draw [rotate=90, fill=white] (0.06,-0.06) circle (0.02);
\end{tikzpicture}
}

\newcommand{\CounterReset}{
\setcounter{CritCounter}{1}
\setcounter{TestCounter}{1}
\setcounter{ColorCounter}{1}
\setcounter{Index}{1}
}

\newcommand\CritLineInitialize[1]{
\ifdef{\SymFun}
{
	\setcounter{CritLineSlide}{#1}
}
{
	\newarray\TestPts
	\newarray\TestVals
	\newarray\TestEVal
	\newarray\fIncDec
	\newarray\testsigns
	\newarray\inequality
	\setcounter{CritLineSlide}{#1}
	\newcommand\SymFun{}
	\newcommand\IncDecStr{}
	\newcommand\TestSignStr{}
	\newcommand\InEqStr{}
	\newcommand\SymFunD{}
	\newcommand\SymFunDLt{}
	\newcommand\SymFunN{}
	\newcommand\SymFunNLt{}
	\newcommand{\CritLineTable}{}
}
}

\newcommand\USubInitialize{
	\newcommand{\Fun}{}
	\newcommand{\FunN}{}
	\newcommand{\FunD}{}
	\newcommand{\FunU}{}
	\newcommand{\UIntegrand}{}
	\newcommand{\AFunU}{}
	\newcommand{\DX}{}
}

\newcommand*\circled[1]{\tikz[baseline=(char.base)]{
            \node[shape=circle,draw,inner sep=1pt,color=blue] (char) {#1};}}

\newcommand*\ellipsed[1]{\tikz[baseline=(char.base)]{
            \node[shape=ellipse,draw,inner sep=1pt,color=blue] (char) {#1};}}

\newcommand*\rected[1]{\tikz[baseline=(char.base)]{
            \node[shape=rectangle,draw,inner sep=1pt,color=blue] (char) {#1};}}

\newcommand*\starred[1]{\tikz[baseline=(char.base)]{
            \node[shape=star,star points=10,draw,inner sep=1pt,color=blue] (char) {#1};}}

\newcommand\CHECKBOX[1]{\makebox[0pt][l]{$\square$}\raisebox{.15ex}{\hspace{0.1em}\ON#1{\B{$\checkmark$}}}}

\makeatletter
\newcommand{\@RemovePositiveSignIfAny}[1]{\IfBeginWith{#1}{+}{\StrGobbleLeft{#1}{1}}{#1}}%
\newcommand{\@PlusIfNonZero}[1]{\IfStrEq{#1}{}{}{+}}%
\newlength{\@MySystemWidth}%
\newcommand*{\@ChooseSignWithAlignment}[2]{%
    \IfBeginWith{#2}{-}%
    {{}-\makebox[\@MySystemWidth][#1]{\ensuremath{\StrGobbleLeft{#2}{1}}}}%
    {{}\@PlusIfNonZero{#2}\makebox[\@MySystemWidth][#1]{\ensuremath{\@RemovePositiveSignIfAny{#2}}}}%
}%
\newcommand*{\@ChooseSignWithRightAlignment}[1]{\@ChooseSignWithAlignment{r}{#1}}%
\newcommand*{\@ChooseSignWithLeftAlignment}[1]{\@ChooseSignWithAlignment{l}{#1}}%

\newenvironment{MySystem}[3][1.2em]{% [<width>] {num of variables} {alignment}
    \newcolumntype{R}{>{\collectcell\@ChooseSignWithRightAlignment}{r}<{\endcollectcell}}
    \newcolumntype{L}{>{\collectcell\@ChooseSignWithLeftAlignment}{l}<{\endcollectcell}}
    \setlength{\@MySystemWidth}{#1}% Store the width to use to typeset the terms
    \IfStrEqCase{#3}{%
        {R}{\begin{array}{r@{} *{\numexpr#2-1\relax}{#3@{}} @{{}={}} r}}%
        {L}{\begin{array}{l@{} *{\numexpr#2-1\relax}{#3@{}} @{{}={}} l}}%
    }[\parbox{\linewidth}{\color{red}MySystem Error: Only L and R alignment types are supported. (using R)}\begin{array}{r *{\numexpr#2-1\relax}{#3} @{{}={}} r}]%
}{%
    \end{array}%
}
\newenvironment{MySystemNoEqual}[3][1.2em]{% [<width>] {num of variables} {alignment}
    \newcolumntype{R}{>{\collectcell\@ChooseSignWithRightAlignment}{r}<{\endcollectcell}}
    \newcolumntype{L}{>{\collectcell\@ChooseSignWithLeftAlignment}{l}<{\endcollectcell}}
    \setlength{\@MySystemWidth}{#1}% Store the width to use to typeset the terms
    \IfStrEqCase{#3}{%
        {R}{\begin{array}{r@{} *{\numexpr#2\relax}{#3@{}} @{{}{}} r}}%
        {L}{\begin{array}{l@{} *{\numexpr#2\relax}{#3@{}} @{{}{}} l}}%
    }[\parbox{\linewidth}{\color{red}MySystem Error: Only L and R alignment types are supported. (using R)}\begin{array}{r *{\numexpr#2\relax}{#3} @{{}{}} r}]%
}{%
    \end{array}%
}

\makeatother


$\quad$

\vspace{0.6cm}

\begin{center}
\LARGE {MA 301} \\

\Large{Final} \\

\large{Spring - 2014} \\

\vspace{0.5cm}

{\setlength{\extrarowheight}{0.2cm}
\begin{tabular}{|c|c|c|}
\hline
Problem & points & out of  \\
\hline
\hline
1 & & 4 \\
\hline
2 & & 4 \\
\hline
3 & & 6 \\
\hline
4 & & 14 \\
\hline
5 & & 6 \\
\hline
6 & & 6 \\
\hline
7 & & 9 \\
\hline
8 & & 15 \\
\hline
9 & & 6 \\
\hline
10 & & 20 \\
\hline
\hline
Score & & 90\\
\hline
\end{tabular}}
\vspace{0.5cm}

%--------TABLE---------
	{\setlength{\extrarowheight}{0.2cm}
	\begin{tabular}{rcc}
		Name: & \underline{$\qquad \qquad \qquad \qquad \quad$}& \\
		&&$\hspace{1.4cm}$\\
		HR: & \underline{$\qquad \qquad \qquad \qquad \quad$}& \\
		&&\\
		Signature: & \underline{$\qquad \qquad \qquad \qquad \quad$}& \\
	\end{tabular}}
%--------------------

\vspace{1.0cm}

	{\setlength{\extrarowheight}{0.2cm}
	\begin{tabular}{| cp{16cm} |}
		\hline
		\multicolumn{2}{|c|}{{\bf READ BEFORE YOU START}} \\
		\hline
		 * & Show all work.\\
		 * & Change your calculator mode to RADIANS.\\
		 * & Sloppy solutions may be penalized.\\
		 * & Use calculators ONLY for computations.\\
		 * & You must write up every problem in a clear and organized way.\\
		 * & If you write solutions on a page different than the question, then clearly state where it is located.\\
		\hline
	\end{tabular}}

\end{center}

\vspace{1.0cm}

\thispagestyle{empty}
\newpage
\thispagestyle{empty}

{\bf Problem 1:} Suppose the force field ${\bf F} = \left[2xy+z, x^2+z^2, x+2yz\right]$ acts upon a point mass along  

some curve $C$ from $A:(1,3,0)$ to $\D{B:\left(1,-1,1\right)}$.

\vspace{0.2cm}

There are two ways to compute the work done in moving the point mass from $A$ to $B$.  Give one of them.

\thispagestyle{empty}
\newpage
\thispagestyle{empty}

{\bf Problem 2:} Give the other method from {\bf Problem 1}.

\thispagestyle{empty}
\newpage
\thispagestyle{empty}

{\bf Problem 3:} Let $\text{\bf F} = [2x^2y,-2xy^2]$ and $C$ be given by $\D{x^2+y^2=5}$ \, oriented counter-clockwise.

Compute \, ${\oint \limits_{C} \text{\bf F}\cdot d\text{\bf r}}$ \, using green's theorem.\\

\thispagestyle{empty}
\newpage
\thispagestyle{empty}

{\bf Problem 4:} Let \,\,$z_1 = -2+2i$,\,\, $z_2 = -i$,\,\, $z_3 = 1-i$. \quad Answer the following in the ({\em specified form}):

\vspace{0.2cm}

(a) ({\em Standard Form}) $z_1z_2 = $

\vspace{1.8cm}

(b) ({\em Standard Form}) $z_3-z_1z_2 = $

\vspace{1.8cm}

(c) $|z_1z_2| = $

\vspace{2.8cm}

(d) Convert $z_2$ to Euler's Polar form:

\vspace{2.8cm}

(e) Convert $z_3$ to Euler's Polar form:

\vspace{2.8cm}

(f) ({\em Euler's Polar Form})  $(z_3)^{76} = $

\vspace{2.8cm}

(g) ({\em Euler's Polar Form}) Solve for $x$: \quad $\D{\frac{1}{\sqrt{2}}x^3 - z_3 = 0}$ \quad (explicitly write out all solutions)

\thispagestyle{empty}
\newpage
\thispagestyle{empty}

{\bf Problem 5:} Let \quad
$
f(x) = 
\left\{
\begin{array}{lrc}
	 \pi & & -\frac{\pi}{2} < x < 0\\
	-\pi & & 0  < x < \frac{\pi}{2}\\
\end{array}
\right.
$  \quad and \quad $f(x) = f(x+\pi)$

\vspace{0.1cm}

Find the Fourier Series of $f(x)$ in terms of: \,  (i) $e^{i\pi n x/L}$ and (ii) $\cos\left(\frac{\pi}{L}nx\right) \, \& \, \sin\left(\frac{\pi}{L}nx\right)$

\thispagestyle{empty}
\newpage
\thispagestyle{empty}

{\bf Problem 6:} \UL{\em Completely Set-up} (simplify and utilize any odd/even properties) the Fourier coefficients for the 

function shown below.  Assume it periodically extends along the $x$-axis.

\hspace{12.0cm} 
\begin{tabular}{cl}
	    & \includegraphics[scale = 0.3]{../../../Figures/Test_Fig19.png} \\
\end{tabular}

\thispagestyle{empty}
\newpage
\thispagestyle{empty}

{\bf Problem 7:} Answer the following true/false and short-answer questions:

\vspace{0.3cm}

{{\setlength{\extrarowheight}{-0.4cm}
\begin{tabular}{|c|c||p{17.2cm}|}
	\hline
	\multirow{6}{*}{T} & \multirow{6}{*}{F} &\\
	&& The line integral of any vector field {\bf F} around any closed curve $C$ is zero\\
	&&\\
	\hline
	\multirow{9}{*}{T} & \multirow{9}{*}{F} &\\
	&& Let $C$ be an ellipse. If $\WORK = 0$, then ${\bf F}$ is path-independent.\\
	&&\\
	\hline
	\multirow{6}{*}{T} & \multirow{6}{*}{F} &\\
	&& If ${\bf F}$ is path-independent, then there is a potential function for ${\bf F}$\\
	&&\\
	\hline
%	\multirow{9}{*}{T} & \multirow{9}{*}{F} &\\
%	&& If $\WORK \ne 0$ for some closed path $C$, then ${\bf F}$ is path-dependent.\\
%	&&\\
%	\hline
	\multirow{9}{*}{T} & \multirow{9}{*}{F} &\\
	&& If ${\bf F}$ is path-independent and $C$ is any closed curve, then $\WORK = 0$\\
	&&\\
	\hline
%	\multirow{10}{*}{T} & \multirow{10}{*}{F} &\\
%	&& If $S_1$ is a rectangle with area = 1 and $S_2$ is a rectangle with area = 2, then
%$\D{2\iint \limits_{S_1} {\bf F} \cdot d{\bf A} = \iint \limits_{S_2} {\bf F} \cdot d{\bf A}}$\\
%	&&\\
%	\hline
	\multirow{10}{*}{T} & \multirow{10}{*}{F} &\\
	&& Suppose $\text{\text{\bf F}} = [1,1,1]$ and $S$ is a closed surface with positive orientation.  Then \,\,$\D{\iint \limits_{S} \text{\bf F}\cdot \text{\bf n} \, dA > 0}$.\\
	&&\\
	\hline
	\multirow{10}{*}{T} & \multirow{10}{*}{F} &\\
	&& Let $S$ be a sphere of radius 1 centered at the origin and oriented outward.  If $\FLUX = 0$, then ${\bf F} = [0,0,0]$\\
	&&\\
	\hline
	\multirow{10}{*}{T} & \multirow{10}{*}{F} &\\
	&& Let $S$ be the surface of a cube centered at the origin with side length 2 and oriented outward.  If ${\bf F} = 2\VI - \VK$ then the flux of ${\bf F}$ through $S$ is 27.\\
	&&\\
	\hline
	\multirow{12}{*}{T} & \multirow{12}{*}{F} &\\
	&& If ${\bf F}$ is a vector-field such that div{\bf F} = 3 and $S$ is a closed surface oriented outward, then $\FLUX$ is equal to one-third the volume enclosed by $S$.\\
	&&\\
	\hline
	\multirow{6}{*}{T} & \multirow{6}{*}{F} &\\
	&& Suppose $f(x) = f(x+6)$ for all $x$, then we say $f$ is periodic with period 12.\\
	&&\\
	\hline
%	\multirow{6}{*}{T} & \multirow{6}{*}{F} &\\
%	&& $\sqrt{-5}\cdot\sqrt{-5} = 5$.\\
%	&&\\
%	\hline
%	\multirow{10}{*}{T} & \multirow{10}{*}{F} &\\
%	&& Let $S$ be a sphere of radius 1 and oriented outward.  If ${\bf F} = [x,y,z]$ then $\FLUX$ is positive\\
%	&&\\
%	\hline
%	\multirow{10}{*}{T} & \multirow{10}{*}{F} &\\
%	&& If $\FLUX = 12$ and $S$ is a flat disk of area $4\pi$, then div${\bf F} = \frac{3}{\pi}$\\
%	&&\\
%	\hline
%	\multirow{9}{*}{T} & \multirow{9}{*}{F} &\\
%	&& If $\WORK = 0$ for one particular closed path, the ${\bf F}$ is path-independent.\\
%	&&\\
%	\hline
%	\multirow{6}{*}{T} & \multirow{6}{*}{F} &\\
%	&& If ${\bf F}$ is path-independent, then there is a potential function for ${\bf F}$\\
%	&&\\
%	\hline
%	\multirow{9}{*}{T} & \multirow{9}{*}{F} &\\
%	&& If $\WORK \ne 0$ for some closed path $C$, then ${\bf F}$ is path-dependent.\\
%	&&\\
%	\hline
%	\multirow{9}{*}{T} & \multirow{9}{*}{F} &\\
%	&& Let $S$ be a sphere of radius 1 centered at the origin and oriented outward.  If ${\bf F} = [x,y,z]$ then $\FLUX$ is positive\\
%	&&\\
%	\hline
%	\multirow{9}{*}{T} & \multirow{9}{*}{F} &\\
%	&& Let $S$ be a sphere of radius 1 centered at the origin and oriented outward.  If $\FLUX = 0$, then ${\bf F} = {\bf 0}$\\
%	&&\\
%	\hline
%	\multirow{9}{*}{T} & \multirow{9}{*}{F} &\\
%	&& Let $S$ be the surface of a cube centered at the origin with side length 2 and oriented outward.  If ${\bf F} = 2\VI - \VK$ then the flux of ${\bf F}$ through $S$ is 27.\\
%	&&\\
%	\hline
%	\multirow{6}{*}{T} & \multirow{6}{*}{F} &\\
%	&& If ${\bf F}$ and ${\bf G}$ are vector-fields satisfying div${\bf F}$ = div${\bf G}$, then ${\bf F} = {\bf G}$\\
%	&&\\
%	\hline
%	\multirow{9}{*}{T} & \multirow{9}{*}{F} &\\
%	&& If ${\bf F}$ is a vector-field such that div{\bf F} = 1 and $S$ is a closed surface oriented outward, then $\FLUX$ is equal to the volume enclosed by $S$.\\
%	&&\\
%	\hline
\end{tabular}}

\thispagestyle{empty}
\newpage
\thispagestyle{empty}

{\bf Problem 8:} Answer the following short-answer questions:

%\setlength{\extrarowheight}{0.2cm}
\begin{tabular}{| p{19cm} |}
	\hline
	\begin{tabular}{p{12cm}c}
		\multicolumn{2}{l}{(a) Let $\D{{\bf F}(x,y,z) = [zy-3y,xz-z,xy+1]}$} \\
		& \multirow{2}{*}{\includegraphics[scale = 0.21]{../../../Figures/10_8_Fig8.png}} \\
		\\ \\ \\ \\ 
		$\D{\int \limits_{C} {\bf F} \cdot d{\bf r} = }$ \\
		\\
	\end{tabular} \\
	\\
	\hline
	(b) Let $\text{\bf F}$ vector field such that curl\,$\text{\bf F} =$ {\bf 0}.
Match each expression in the left column with the most appropriate term in the right column.
	\begin{center}
	%{{\setlength{\extrarowheight}{0.1cm}
	\begin{tabular}{cp{2cm}c}
		$\D{\int \limits_{C} \text{\bf F}\cdot d\text{\bf r}}$ & & conservative\\
		$\D{\oint \limits_{D} \text{\bf F}\cdot d\text{\bf r}}$ & & path independent\\
		$\D{\text{\bf F}}$ & & zero\\
	\end{tabular}
	%}}
	\end{center}
	\\
\hline
	(c) Match each the Theorem in the left column with the most appropriate statement in the right column.
	\begin{center}
	{\setlength{\extrarowheight}{0.3cm}
	\begin{tabular}{rp{1.6cm}l}
		Divergence Theorem & & The surface of a full sphere of radius 4\\
		Stokes' Theorem & & The surface of a half-sphere of radius 4\\
		Green's Theorem & & A disk of radius 4\\
	\end{tabular}}
	\end{center}
	\\
	\hline
	\begin{tabular}{p{12cm}c}
		\multicolumn{2}{l}{(d) Let $\D{{\bf F} = \nabla\left(\frac{y-1}{y+z+1}\right)}$ and let\, $C$ and $D$ be shown in the figure. Then} \\
		& \multirow{2}{*}{\includegraphics[scale = 0.25]{../../../Figures/10_8_Fig6.png}} \\
		\\
		$\D{\int \limits_{C} {\bf F} \cdot d{\bf r} = }$ \\
		&\\
		\\
		\\
		$\D{\int \limits_{D} {\bf F} \cdot d{\bf r} = }$ &\\
		\\ \\ \\
	\end{tabular}
	\\ \\
	\hline
	(e) Let $\text{\bf F} = [6x,ye^z,-e^z]$ and $S$ be the surface (oriented outward) of the box $|x|\le 1$, $|y|\le 2$ , $0\le z \le 1$.
	\\ \\ \\ \\ \\
	$\D{\iint \limits_{S} \text{\bf F}\cdot \text{\bf n} \, dA =}$
	\\ \\ \\ \\ \\
	\hline
\end{tabular}

\thispagestyle{empty}
\newpage
\thispagestyle{empty}

{\bf Problem 9:} \UL{\em Set-up the integrals} (to the point you would use an integration technique) needed to compute:

(a) The mass of the wire with shape given by \RT = t\VI + 2t\VJ + $\frac{2}{3}t^{3/2}$\VK, \, $0 \le t \le 2$ \& density $\delta(x,y,z) = 3\sqrt{5+x}$.

\vspace{11.0cm}

(b) The surface area of $S$ where $S$ is given by the parameterization: ${\bf r}(u,v) = [u,v,u^3]$,\,\, $0 \le u \le 1$,\,\, $-2 \le v \le 2$

\vspace{8.0cm}

%(c) The flux of ${\bf F} = [x^2,0,z^2]$ \, through the closed rectangular box: $|x| \le 1$, \,\, $|y| \le 3$, \,\, and $0 \le z \le 2$.
%
%\vspace{0.2cm}

\thispagestyle{empty}
\newpage
\thispagestyle{empty}

{\bf Problem 10:} Answer all of the following questions based on the figure below.

{\setlength{\extrarowheight}{-0.3cm}
\begin{tabular}{p{0.35cm}p{12cm}p{4cm}}
	&& \multirow{4}{*}{\includegraphics[scale = 0.32]{../../../Figures/Test_Fig16.png}} \\
	\multicolumn{2}{c}{$\D{{\bf F}(x,y,z) = \left[x+y,y^2,4-z\right]}$} & \\
	&&\\
	&&\\
	$S_B$: & Part of the $z=0$ plane with $-2 \le x \le 2$, $0\le y \le 3$ & \\
%	&&\\
%	$C_B$: & Boundary of $S_B$ & \\
	&&\\
	$S_L$: & Part of the $y=0$ plane with $0 \le z \le 4-x^2$, $-2 \le x \le 2$ & \\
%	&&\\
%	$C_L$: & Boundary of $S_L$ & \\
	&&\\
	$S_R$: & Part of the $y=3$ plane with $0 \le z \le 4-x^2$, $-2 \le x \le 2$ & \\
%	&&\\
%	$C_R$: & Boundary of $S_R$ & \\
	&&\\
	$S_T$: & Part of the parabolic cylinder $z=4-x^2$ with $0\le y \le 3$, $-2 \le x \le 2$ & \\
%	&&\\
%	$C_T$: & Boundary of $S_T$ & \\
	&&\\
	$S$: & Total closed surface in the figure (i.e. $S = S_B \cup S_L \cup S_R \cup S_T$) & \\
\end{tabular}}

\vspace{0.2cm}

{\em Hint}: ${\text{Area}(S_L) = \text{Area}(S_R) = 2\int \limits_{0}^{2} (4-x^2) \, dx}$
\vspace{0.2cm}

{\setlength{\extrarowheight}{0.1cm}
\begin{tabular}{p{0.35cm}p{20.6cm}}
	(a) & Compute the flux through $S_B$ directly.\\
	& \hspace{12.5cm} ({\em Assume outward orientation})\\
	&\\&\\&\\&\\&\\&\\&\\&\\&\\&\\
	(b) & Compute the flux through $S_L$ directly. \\
	&\\&\\&\\&\\&\\&\\&\\&\\&\\&\\&\\
	(c) & Compute the flux through $S_R$ directly. \\
	&\\
\end{tabular}}

\thispagestyle{empty}
\newpage
\thispagestyle{empty}

{\setlength{\extrarowheight}{0.1cm}
\begin{tabular}{p{0.35cm}p{14.6cm}}
	(d) & Compute the flux through $S_T$ directly. \\
	&\\&\\&\\&\\&\\&\\&\\&\\&\\&\\&\\&\\&\\&\\&\\&\\&\\&\\&\\&\\&\\&\\&\\&\\&\\&\\&\\&\\&\\&\\&\\&\\&\\&\\&\\&\\&\\
	(e) & Compute the flux through $S$ using (a)-(d). \\
	&\\
\end{tabular}}

\thispagestyle{empty}
\newpage
\thispagestyle{empty}

{\setlength{\extrarowheight}{0.1cm}
\begin{tabular}{p{0.35cm}p{18.6cm}}
	(f) & Compute the flux through $S$ using the Divergence Theorem. \\
	&\\&\\&\\&\\&\\&\\&\\&\\&\\&\\&\\&\\&\\&\\&\\&\\&\\&\\&\\&\\&\\&\\&\\&\\&\\&\\&\\
	(g) & If $C_T$ is the boundary of $S_T$ ({\em oriented counter-clockwise}), compute \, $\D{\int \limits_{C_T} {\bf F}\,\cdot\, d{\bf r}}$ \,\, using Stokes' Thm \\
	&\\
\end{tabular}}

%{\bf Problem 10:} Use Stokes' Thm to answer all the questions below.  (You may use any results from {\bf Problem 2}) 
%
%{\setlength{\extrarowheight}{-0.3cm}
%\begin{tabular}{p{0.35cm}p{12cm}p{4cm}}
%	&& \multirow{4}{*}{\includegraphics[scale = 0.32]{../../../Figures/Test_Fig18.png}} \\
%	\multicolumn{2}{c}{$\D{{\bf F}(x,y,z) = \left[x+y,y^2,-4-z\right]}$} & \\
%	&&\\
%	&&\\
%	$S_B$: & Part of the $z=0$ plane with $-2 \le x \le 2$, $0\le y \le 3$ & \\
%	&&\\
%	$S_L$: & Part of the $y=0$ plane with $0 \le z \le 4-x^2$, $-2 \le x \le 2$ & \\
%	&&\\
%	$S_R$: & Part of the $y=3$ plane with $0 \le z \le 4-x^2$, $-2 \le x \le 2$ & \\
%	&&\\
%	$S_T$: & Part of the parabolic cylinder $z=4-x^2$ with $0\le y \le 3$, $-2 \le x \le 2$ & \\
%	&&\\
%\end{tabular}}
%
%\vspace{0.2cm}
%
%{\setlength{\extrarowheight}{0.1cm}
%\begin{tabular}{p{0.35cm}p{11.6cm}}
%	(a) & Compute $\D{\int \limits_{C_B} {\bf F}\,\cdot\, d{\bf r}}$ \quad where $C_B$ is the boundary of $S_B$\\
%	&\\&\\&\\&\\&\\&\\&\\&\\&\\&\\
%	(b) & Compute $\D{\int \limits_{C_L} {\bf F}\,\cdot\, d{\bf r}}$ \quad where $C_L$ is the boundary of $S_L$\\
%	&\\&\\&\\&\\&\\&\\&\\&\\&\\&\\
%	(c) & Compute $\D{\int \limits_{C_T} {\bf F}\,\cdot\, d{\bf r}}$ \quad where $C_T$ is the boundary of $S_T$\\
%	&\\
%\end{tabular}}
%
%\thispagestyle{empty}
%\newpage
%\thispagestyle{empty}

\end{document}